\section{Work Packages and Timeline}

The project is organized into ten dependent work packages over 22 months. Promotion criteria are evidence gates rather than target dates alone.

\begin{table*}[t]
	\centering
	\caption{Implementation work packages.}
	\label{tab:work-packages}
	\begin{tabularx}{\textwidth}{p{0.07\textwidth}p{0.1\textwidth}p{0.24\textwidth}X}
		\toprule
		WP & Months & Scope                                                                 & Acceptance criterion                                         \\
		\midrule
		1  & 1--2   & Scope, inventory boundaries, requirements, ethics and hazard analysis & Signed requirements, control exclusions and governance roles \\
		2  & 1--3   & QIDK/Android qualification and power method                           & Reproducible hardware/runtime report and soak test           \\
		3  & 2--5   & Metering, gateway, Brick graph and quality pipeline                   & At least 98\% valid critical intervals over 14 days          \\
		4  & 3--7   & Energy/carbon baselines and M\&V plan                                 & Frozen model passing predeclared residual diagnostics        \\
		5  & 4--8   & Forecasting, anomaly detection, uncertainty and quantization          & Candidate beats declared baselines and QIDK release gates    \\
		6  & 6--10  & Digital twin, MPC and independent safety kernel                       & BOPTEST and HIL fault suite pass                             \\
		7  & 9--12  & Live shadow and advisory deployment                                   & Stable availability; no uncontrolled writes                  \\
		8  & 12--18 & Registered randomized or crossover field trial                        & Sample, data quality and safety criteria met                 \\
		9  & 16--20 & Transfer, ablation, privacy and resilience studies                    & Locked cross-building and backend comparisons                \\
		10 & 19--22 & M\&V, audit, artifacts and dissemination                              & Independent reproducibility review and final release         \\
		\bottomrule
	\end{tabularx}
\end{table*}

Major milestones are: M1 scope approval; M2 privacy, security and safety review; M3 QIDK qualification; M5 telemetry commissioning; M7 frozen baselines; M9 HIL completion; M10 live shadow; M12 guarded-control readiness; M18 field-trial completion; M20 locked ledger; M21 reproducibility audit; and M22 final report and release.

Deliverables are: a boundary and governance specification; QIDK benchmark report; Brick point graph; versioned telemetry pipeline; baseline and predictive models; constrained optimizer and safety kernel; signed Android package; HIL suite; registered analysis plan; privacy-preserving research dataset; M\&V and net-zero ledger; model cards, data sheets, threat model and privacy assessment; reproducibility bundle; and final manuscript.

\subsection{Resource Plan}

Minimum roles are a principal investigator, building-controls engineer, Android/edge engineer, ML researcher, energy/M\&V analyst, data engineer, cybersecurity reviewer, privacy/ethics reviewer, and facilities operator. The pilot requires the QIDK target, managed Android support, industrial protocol gateway, meter access, calibrated device-level power analyzer, development server, test BMS or HIL interface, and approved access to weather and grid-factor data.

Procurement must include meter calibration and installation, gateway support, environmental sensors, secure mounting, power instrumentation, and software or standards access. The embodied impact and operational consumption of added equipment are recorded so the project does not externalize its own footprint.
